\documentclass[11pt,a4paper]{article}

\usepackage[margin=1in]{geometry}
\usepackage{titlesec}
\usepackage{enumitem}
\usepackage{hyperref}
\usepackage{parskip}
\usepackage{graphicx}
\usepackage{float}
\usepackage{booktabs}
\usepackage{amsmath}

\hypersetup{colorlinks=true, linkcolor=blue, urlcolor=blue}

\titleformat{\section}{\large\bfseries}{\thesection}{0.5em}{}
\titleformat{\subsection}{\normalsize\bfseries}{\thesubsection}{0.5em}{}
\setlist{leftmargin=1.4em, topsep=2pt, itemsep=2pt}

\title{\vspace{-1.5cm}\textbf{PKCERT AI \& Software Development Internship \\ Task 28: Production LLM Serving, Asynchronous Microservices \& Full-Stack Deployment}}
\author{Abdullah Amir}
\date{AI and Software Development \\ 24 August 2026}

\begin{document}

\maketitle
\vspace{-0.8cm}

Full code lives in this task's directory (\texttt{app/}, \texttt{frontend/}, \texttt{Dockerfile},
\texttt{docker-compose.yml}, \texttt{benchmark/}). An asynchronous FastAPI microservice serves
three Hugging Face \texttt{transformers} pipelines (sentiment, summarization, generation),
containerized via a multi-stage Docker build, orchestrated locally via \texttt{docker-compose},
and paired with a self-contained HTML/CSS/JS frontend served from the same container. The
Docker image was built and run as a real container, verified working end to end (non-root,
healthy, all four endpoints responding). Live cloud deployment was attempted on Render and
hit a genuine, measured 512MB memory ceiling -- Part C.3 documents this in full, including two
mitigations that were actually implemented and tested (not merely proposed) and did not
resolve it.

% ----------------------------------------------------------------
\section*{Part A: Pretrained LLM Inference \& Multi-Task Orchestration}

\textbf{A1 -- models, and why.}

\begin{table}[H]
\centering
\small
\begin{tabular}{@{}p{1.7cm}p{4.1cm}rp{6.8cm}@{}}
\toprule
\textbf{Task} & \textbf{Model} & \textbf{Params} & \textbf{Rationale} \\
\midrule
Sentiment & \raggedright cardiffnlp/twitter-\allowbreak roberta-base-\allowbreak sentiment-latest & 125M & Multi-class (neg/neu/pos), not the more common binary SST-2 default -- the brief specifically asks for multi-class. \\[0.4em]
Summarization & \raggedright t5-small & 60M & Deliberately the lightest viable abstractive summarizer, not a BART-family model (306M-400M) -- a direct instance of A4's memory-footprint trade-off. \\[0.4em]
Generation & \raggedright distilgpt2 & 82M & The standard lightweight causal LM for CPU-only serving. \\
\bottomrule
\end{tabular}
\end{table}

Combined resident footprint once all three are loaded: \textbf{$\sim$695MB RSS} (measured
directly, D2) -- a number that turned out to matter far more than expected once real cloud
deployment was attempted (C3).

\textbf{A2 -- decoding strategies, compared with real output.} All five are implemented and
selectable per-request via \texttt{decoding\_strategy}. Prompt: \emph{"The best way to learn
programming is"}, \texttt{max\_new\_tokens=30}:

\begin{table}[H]
\centering
\small
\begin{tabular}{@{}p{2cm}p{6.3cm}p{1.6cm}p{2.5cm}@{}}
\toprule
\textbf{Strategy} & \textbf{Sample output (truncated)} & \textbf{Diversity} & \textbf{Coherence} \\
\midrule
Greedy & "...to learn the basics of programming.\textbackslash n\textbackslash n\textbackslash n..." & None & Degenerates into repetition \\[0.3em]
Beam (4) & "...to learn the basics of programming.\textbackslash n\textbackslash n..." & Low & Same repetition failure mode \\[0.3em]
Top-k (50) & "...to take time off of school to become a better, happier, and smarter person." & Moderate & Good \\[0.3em]
Top-p (0.9) & "...to develop your own languages. We've seen how programming languages do..." & High & Good \\[0.3em]
Temperature (1.0) & "...to set aside your mind to do so. The best way to teach this is..." & Highest & Lower \\
\bottomrule
\end{tabular}
\end{table}

Greedy and beam search both degenerated into repetition on this open-ended prompt -- a real,
observed failure mode, not just textbook description: both always pick the (locally or
globally) highest-probability continuation, which for free-form generation tends toward
short repeating loops once the model enters one. Top-k, top-p, and temperature sampling all
avoided this by introducing randomness into token selection, at the cost of losing greedy's
determinism. \textbf{Latency impact}: decoding strategy affects output quality far more than
per-token latency at this model size -- all five perform one forward pass per generated
token; beam search is the one real exception, since it evaluates \texttt{num\_beams}
candidate sequences per step rather than one.

\textbf{A3 -- payload sanitization \& context-window management.} Input is stripped of
non-printable characters, tokenized with the \emph{target model's own tokenizer} (word count
and subword-token count diverge -- Task 27's finding), and truncated at the token level (never
mid-token) when it exceeds a per-task limit, rather than erroring outright. \textbf{A genuine
bug this surfaced}: naively truncating token ids, decoding back to text, and handing that text
to the pipeline is not token-count-lossless -- the pipeline re-tokenizes the decoded text from
scratch and adds fresh special tokens, which can push the count back over a model's true limit.
This produced a real \texttt{IndexError: index out of range in self} from RoBERTa's
position-embedding lookup on an 800-word test input. Fixed with a bounded re-check-and-re-trim
loop (D3 has the full account); re-verified afterward with the same input, now truncating
cleanly to 511 tokens with no error.

\textbf{A4 -- architectural trade-off: local weights vs.\ remote inference API.}

\begin{table}[H]
\centering
\small
\begin{tabular}{lp{5.6cm}p{5.6cm}}
\toprule
\textbf{Dimension} & \textbf{Direct weight loading (this task)} & \textbf{Remote inference API} \\
\midrule
Memory & Full model resident in-process ($\sim$695MB, all three) & Zero local model memory \\
Compute & Every request pays full CPU forward-pass cost & Runs on provider's (often GPU) infra \\
Latency & No extra network hop; 24--500ms once warm & Extra round trip, but often GPU-fast \\
Cold start & Measured, three tiers: 171s (download), 2.2--6.9s (cache-warm), sub-second (fully warm) & Opaque, provider-controlled \\
Rate limits & None -- bounded only by this host's CPU/RAM & Provider-imposed, often stricter on free tiers \\
\bottomrule
\end{tabular}
\end{table}

Direct loading wins here specifically because the three models are small (267M params
combined) -- a 70B-parameter model would make direct CPU loading impractical on memory alone,
at which point a remote GPU-backed API becomes the only viable option regardless of its
network/rate-limit costs.

% ----------------------------------------------------------------
\section*{Part B: Asynchronous REST API Engineering}

\textbf{B1 -- endpoints.} \texttt{POST /api/v1/sentiment}, \texttt{/summarize},
\texttt{/generate}, \texttt{GET /healthz}.

\textbf{B2 -- validation \& error handling.} Pydantic v2 models enforce type/length/range
bounds (\texttt{temperature} $\in[0.01,2.0]$, \texttt{top\_p}$\in[0,1]$,
\texttt{max\_new\_tokens}$\in[1,256]$) \emph{before} any request reaches inference code --
verified directly: an out-of-range temperature (5.0) returned \texttt{422} with a structured
body before touching a model. Dedicated handlers map \texttt{ContextWindowError}$\to$400,
\texttt{RequestValidationError}$\to$422, and any unhandled exception$\to$500, all sharing one
\texttt{\{error, detail, status\_code\}} response shape.

\textbf{B3 -- concurrency \& state.} Every Hugging Face pipeline call is synchronous, CPU-bound
Python -- calling it directly inside an \texttt{async def} route would still block the single
event loop for the full inference duration. Every inference call is dispatched via
\texttt{run\_in\_threadpool} instead, keeping the event loop free to route other requests
(including \texttt{/healthz}) while CPU-bound work runs on a worker thread. \texttt{get\_pipeline()}
is a lock-guarded lazy-loading singleton: the first request for a task pays its load cost once;
every later request (any caller) reuses the cached pipeline. \texttt{PRELOAD\_MODELS=true}
(the Docker default) instead pays all three models' cost once at startup -- A4's cold-start
table is the direct evidence behind this trade-off. Structured JSON request logging and CORS
middleware are both configured.

% ----------------------------------------------------------------
\section*{Part C: Containerization \& Cloud Infrastructure Deployment}

\textbf{C1 -- multi-stage Dockerfile.} Stage \texttt{builder} (python:3.11-slim) resolves and
installs dependencies -- including the heavy \texttt{torch}/\texttt{transformers} wheels --
into an isolated virtualenv, in its own layer; since \texttt{requirements.txt} is copied and
installed \emph{before} application code, editing \texttt{app/*.py} never invalidates this
layer, so a code-only rebuild skips re-downloading torch entirely. Stage \texttt{runtime}
starts from a fresh minimal base and copies over only the built virtualenv + source -- none of
the builder stage's pip cache or intermediate layers survive into the final image. Runs as a
dedicated non-root user; includes a container-level \texttt{HEALTHCHECK}. \textbf{Measured
final image size: 1.73GB} -- critically, this required explicitly pinning
\texttt{pip install torch --index-url https://download.pytorch.org/whl/cpu} \emph{before} the
general \texttt{requirements.txt} install; a plain \texttt{pip install torch} instead pulls
PyPI's default CUDA-enabled build, bundling multi-hundred-MB NVIDIA driver libraries this
CPU-only deployment target never uses -- caught by watching the first build's download log
show an unexpectedly large torch wheel mid-build, not assumed in advance. \textbf{Verified end
to end, not just built}: \texttt{docker run} on the resulting image loaded all three models
(matching A4's cold-start timings), \texttt{docker exec ... whoami} confirmed the container
runs as the non-root \texttt{app} user, Docker's own \texttt{HEALTHCHECK} reported
\texttt{healthy}, and all three API endpoints plus the static frontend responded correctly
against the live container.

\textbf{C2 -- local orchestration.} \texttt{docker-compose.yml} configures environment variable
injection (\texttt{PORT}, \texttt{PRELOAD\_MODELS}, optional \texttt{.env} for
\texttt{HF\_TOKEN}), an automated health check (240s start period, since preloading pays all
three models' cost before reporting healthy), host port 8000 mapped to container port 7860, a
named volume persisting the Hugging Face model cache across restarts, and explicit CPU
(2.0)/RAM (2GB limit, 1GB reservation) quotas.

\textbf{C3 -- cloud deployment: attempted, with a genuine resource-ceiling finding.}
\textbf{Hugging Face Spaces was the original target}, but a direct API call
(\texttt{POST /api/repos/create} for a Docker-SDK Space) returned \texttt{402 Payment
Required} -- Hugging Face now requires a paid PRO subscription to host Docker/Gradio Spaces
even on the free tier; only static Spaces (no backend) remain free. Pivoted to
\textbf{Render} (also brief-approved).

The service was created and deployed via Render's own API (Docker runtime, root directory
\texttt{Day\_28}, auto-deploy on push to \texttt{main} -- Render's native commit-triggered
rebuild \emph{is} the CI/CD pipeline Part C.3 asks for, no separate GitHub Actions workflow
needed). The build succeeded in $\sim$2.5 minutes. The \textbf{runtime deploy failed}:
\texttt{Out of memory (used over 512Mi)} -- Render's free tier caps a service at 512MB, and
this service needs $\sim$695MB with all three models resident (D2). Render's logs API showed
the OOM was hit loading the \emph{first} model, before the other two even began.

\textbf{Two mitigations were implemented and measured, not merely proposed}: (1)
\texttt{torch\_dtype=\allowbreak torch.bfloat16} (halving the theoretical per-parameter cost) --
measured RSS after loading all three: \textbf{998MB}, \emph{higher} than the fp32 baseline
(695MB), likely because this environment's CPU inference path upcasts back to fp32 for some
ops. (2) \texttt{torch.quantization.quantize\_dynamic} (int8, applied to the sentiment model
as an isolated test) -- measured RSS: \textbf{1024MB}, worse still, since the (now itself
deprecated) API produces a new quantized model while the original fp32 model remains
resident unless manually freed. Both results are reported as measured, precisely because an
``obvious'' fix that doesn't hold up under direct measurement is a more honest finding than
one assumed to work.

\textbf{Local Docker verification} (fully confirmed, separate from the cloud attempt): image
built (\texttt{docker build}, 1.73GB final size, no bloat from unwanted CUDA libraries --
Part C.1's index-URL fix), container run (\texttt{docker run}), non-root user confirmed
active inside the running container, and all four endpoints exercised directly against the
live container. One further genuine finding surfaced here: the container initially could not
resolve DNS at all inside this development sandbox's default Docker bridge network (a
\texttt{NameResolutionError} on every Hugging Face Hub request, despite the host machine's
own unrestricted internet access) -- resolved by passing explicit DNS resolvers
(\texttt{--dns=8.8.8.8 --dns=1.1.1.1}, now also in \texttt{docker-compose.yml}), after which
the container downloaded and served all three models correctly.

The Render service (\texttt{task28-llm-microservice}) exists and is configured correctly but
is not currently serving traffic. Genuinely fitting this architecture into 512MB would need
either a paid tier (Render's Starter plan, $\sim$\$7/month, raises the limit to 2GB) or a
materially different approach than what was tried (ONNX Runtime's own quantization tooling,
or splitting the three models across separate lightweight services) -- documented here as the
concrete next step rather than silently worked around.

% ----------------------------------------------------------------
\section*{Part D: Frontend, Load Testing \& Documentation}

\textbf{D1 -- frontend.} Plain HTML/CSS/JS, no framework, served by FastAPI itself from the
same container -- one deployable unit. Three tabs, live parameter controls for generation
(the strategy dropdown shows/hides only the relevant sliders), a polling health badge, and
per-request status/error reporting. Every request uses relative URLs, so the identical static
files work unmodified against localhost, a local container, or any future live cloud URL
(untested against a live cloud instance -- see C3). Verified
in an actual browser (headless Chromium screenshot) during development -- the health badge
correctly reported "API online $\cdot$ 44ms $\cdot$ 3 model(s) loaded" once all three
pipelines were warm.

\textbf{D2 -- load testing.} Single-client (sequential) vs.\ 6-way concurrent, against the
local instance with all three models pre-warmed:

\begin{table}[H]
\centering
\small
\begin{tabular}{llrrrr}
\toprule
\textbf{Endpoint} & \textbf{Pattern} & \textbf{P50} & \textbf{P90} & \textbf{P99} & \textbf{Throughput} \\
\midrule
sentiment & sequential & 24.9ms & 28.6ms & 37.0ms & 34.3 req/s \\
sentiment & 6 concurrent & 81.2ms & 84.3ms & 86.1ms & 69.5 req/s \\
summarize & sequential & 491.8ms & 520.5ms & 541.8ms & 2.02 req/s \\
summarize & 6 concurrent & \textbf{3997.8ms} & 4058.6ms & 4079.3ms & 1.54 req/s \\
generate & sequential & 308.6ms & 343.1ms & 352.0ms & 3.20 req/s \\
generate & 6 concurrent & \textbf{1467.8ms} & 1524.9ms & 1527.7ms & 4.07 req/s \\
\bottomrule
\end{tabular}
\end{table}

Peak RSS: 397MB (sentiment load), 399MB (summarize load), 686MB (generate load, all three
models resident by this point). Full per-run detail in \texttt{benchmark/results/local.json}.
This machine had real background load from other running applications during measurement, so
these are \emph{this-environment} numbers, not clean-room absolutes -- reported transparently
rather than presented as idealized.

\textbf{The real finding}: sentiment (the cheapest model) scales throughput well under
concurrency (34$\to$69 req/s) with latency growing only $\sim$3x for a 6x concurrency
increase. Summarize and generate (the multi-forward-pass generation tasks) show the opposite
pattern -- throughput barely moves (summarize: 2.02$\to$1.54 req/s, actually \emph{worse})
while per-request latency balloons 5--8x. This is direct evidence of a CPU-bound concurrency
bottleneck: \texttt{torch} releases Python's GIL during its C++ tensor ops, so
\texttt{run\_in\_threadpool} genuinely allows some overlap, but six simultaneous multi-step
generation loops still compete for the same limited CPU cores, and the scheduler dividing
those cores six ways serializes the \emph{effective} per-request throughput almost entirely
for the heavier tasks.

\textbf{D3 -- two non-trivial engineering challenges.}

\emph{1. An unpinned dependency install silently pulled a breaking major version.}
\texttt{pip install transformers} installed \texttt{transformers==5.15.1} (the newest release
at build time), which no longer registers \texttt{"summarization"} as a valid pipeline task
(\texttt{KeyError: Unknown task summarization}, with the printed available-tasks list
conspicuously omitting both \texttt{summarization} and \texttt{text2text-\allowbreak generation}
entirely). Diagnosed by reading the actual exception's task list rather than assuming a typo
in this codebase; resolved by pinning to \texttt{transformers==4.57.6}, restoring the expected
pipeline registry -- the same class of "an unpinned major-version jump breaks an established
API" issue this internship already documented once, in Task 25 (gensim vs.\ Python 3.14).
\texttt{requirements.txt} now pins every dependency with \texttt{==}, not a floor, specifically
so this failure mode is reproducible-and-fixable rather than silently different on the next
install.

\emph{2. The concurrency bottleneck measured in D2.} Rather than assume
\texttt{run\_in\_threadpool} was sufficient for real concurrency, D2's load test was run
deliberately, and it revealed that it isn't -- not because the async wiring is wrong, but
because six CPU-bound generation loops on a host with limited cores are fundamentally
compute-bound, and no amount of async/threading changes that. \textbf{The concrete fix this
motivates, matching the brief's own suggested mitigations}: (a) \emph{worker-pool adjustment}
-- running multiple \texttt{uvicorn} worker \emph{processes} (\texttt{--workers N}) would let
the OS schedule genuinely independent processes across cores instead of Python threads sharing
one interpreter, at the cost of N$\times$ the model memory (a real trade-off against A1/A4's
small-model strategy, not a free win); (b) \emph{quantization} --
\texttt{torch.quantization.quantize\_dynamic} on the linear layers would reduce each forward
pass's compute cost directly, lowering both single-request latency and the per-request cost
that compounds under concurrent load, without the memory multiplication worker processes
require. Neither was applied in this session (documented here as the identified, justified
next step rather than implemented against a deadline).

\textbf{D4 -- reflection.} Static-vs-dynamic trade-offs recur throughout this task: small
models chosen over larger ones for deployability (A1), lazy-loading traded against startup
latency (B3/A4), and concurrency itself trading latency for throughput unevenly across the
three task types (D2). No single configuration is correct independent of the deployment
target -- a resource-constrained free tier and a dedicated multi-core production host would
reasonably make different choices at every one of these decision points.

\end{document}
